\documentclass[10pt]{article}
\providecommand{\FIGDIR}{fig}
\newcommand{\DIGESTDATE}{8 August 2026}
\newcommand{\PERIODSTART}{2 August 2026}
\newcommand{\PERIODEND}{8 August 2026}
% ===================== PM Research Watch - shared preamble =====================
\usepackage[T1]{fontenc}
\usepackage[utf8]{inputenc}
\usepackage{amsmath,amssymb}
\usepackage{tgpagella}
\usepackage{tgheros}
\renewcommand{\sfdefault}{qhv}
\renewcommand{\ttdefault}{lmtt}
\usepackage{microtype}
\usepackage[english]{babel}

\usepackage[a4paper,top=17mm,bottom=17mm,left=17mm,right=17mm,
            headsep=6mm,footskip=10mm]{geometry}
\usepackage{graphicx}
\usepackage{xcolor}
\usepackage{colortbl}
\usepackage{tikz}
\usetikzlibrary{calc}
\usepackage[most]{tcolorbox}
\usepackage{titlesec}
\usepackage{fancyhdr}
\usepackage{booktabs}
\usepackage{array}
\usepackage{longtable}
\usepackage{enumitem}
\usepackage{multicol}
\usepackage{ragged2e}
\usepackage[hidelinks]{hyperref}

% ------------------------------------------------------------------ palette
\definecolor{Ink}    {HTML}{14262E}
\definecolor{Deep}   {HTML}{0F4C5C}
\definecolor{Teal}   {HTML}{1F7A8C}
\definecolor{Sky}    {HTML}{4FA3C4}
\definecolor{Sage}   {HTML}{6FA287}
\definecolor{Amber}  {HTML}{E3A008}
\definecolor{Clay}   {HTML}{B5643C}
\definecolor{Coral}  {HTML}{D1495B}
\definecolor{Violet} {HTML}{6A5B9E}
\definecolor{Slate}  {HTML}{4A6572}
\definecolor{Paper}  {HTML}{FBF9F5}
\definecolor{Mist}   {HTML}{EEEAE2}
\definecolor{Rule}   {HTML}{D9D3C9}

\pagecolor{Paper}
\color{Ink}

\hypersetup{colorlinks=true, linkcolor=Deep, urlcolor=Teal, citecolor=Teal,
            pdftitle={PM Research Watch}, pdfauthor={Automated daily digest}}

% ------------------------------------------------------------------ headings
\titleformat{\section}
  {\sffamily\bfseries\large\color{Deep}}{}{0pt}
  {}[{\vspace{1.2mm}\color{Amber}\titlerule[1.1pt]}]
\titlespacing*{\section}{0pt}{5.5mm}{2.8mm}

\titleformat{\subsection}
  {\sffamily\bfseries\normalsize\color{Teal}}{}{0pt}{}
\titlespacing*{\subsection}{0pt}{3.6mm}{1.4mm}

\setlist[itemize]{leftmargin=4.2mm, itemsep=1.1mm, topsep=1.3mm, parsep=0pt,
                  label=\textcolor{Teal}{\footnotesize$\blacktriangleright$}}

% ------------------------------------------------------------------ header/footer
\pagestyle{fancy}
\fancyhf{}
\renewcommand{\headrulewidth}{0pt}
\renewcommand{\footrulewidth}{0pt}
\fancyhead[L]{\sffamily\scriptsize\color{Slate}PM RESEARCH WATCH\;\textcolor{Rule}{\textbar}\;\DIGESTDATE}
\fancyhead[R]{\sffamily\scriptsize\color{Slate}Particulate matter monitoring \& health effects}
\fancyfoot[C]{\sffamily\scriptsize\color{Slate}\thepage}

% ------------------------------------------------------------------ components
% metric tile
\newcommand{\metric}[3]{%
  \begin{tcolorbox}[enhanced, colback=white, colframe=Rule, boxrule=0.5pt,
    arc=1.6mm, left=2.2mm, right=2.2mm, top=1.8mm, bottom=1.8mm,
    borderline west={1.7mm}{0pt}{#3}, nobeforeafter, valign=center,
    height=17mm, width=\linewidth]
    \RaggedRight
    {\sffamily\bfseries\fontsize{17}{17}\selectfont\color{#3}#1}\\[0.5mm]
    {\sffamily\fontsize{6.5}{7.4}\selectfont\color{Slate}#2}
  \end{tcolorbox}}

% pull-quote / key insight box
\newtcolorbox{keybox}[1][]{enhanced, colback=Mist, colframe=Mist,
  arc=1.6mm, left=3.2mm, right=3.2mm, top=2.4mm, bottom=2.4mm,
  borderline west={1.7mm}{0pt}{Amber}, #1}

% subtopic band
\newcommand{\band}[2]{%
  \vspace{3.2mm}\noindent
  \begin{tcolorbox}[enhanced, colback=#2, colframe=#2, arc=1.2mm,
    left=2.6mm, right=2.6mm, top=1.1mm, bottom=1.1mm, boxrule=0pt,
    width=\linewidth, nobeforeafter]
    {\sffamily\bfseries\small\color{white}#1}
  \end{tcolorbox}\vspace{1.6mm}}

% one paper entry: {short cite}{title}{journal}{doi}{body}
\newcommand{\paperentry}[5]{%
  \par\vspace{1.9mm}\noindent
  {\sffamily\bfseries\footnotesize\color{Ink}#1}\;%
  {\sffamily\fontsize{7}{8}\selectfont\color{Slate}#3}\\[0.4mm]
  {\footnotesize\itshape\color{Deep}#2}\\[0.6mm]
  {\footnotesize #4}\\[0.3mm]
  {\fontsize{6.6}{7.6}\selectfont\color{Slate}\texttt{doi:}\,\href{https://doi.org/#5}{\texttt{\detokenize{#5}}}}
  \par}

% caption style
\newcommand{\figcap}[1]{\par\vspace{1.0mm}
  {\fontsize{7.2}{8.6}\selectfont\color{Slate}\textsf{\textbf{Fig.}}\; #1\par}\vspace{1.5mm}}

\providecommand{\permil}{\textperthousand}

\setlength{\parindent}{0pt}
\setlength{\parskip}{0pt}
\renewcommand{\arraystretch}{1.18}


\begin{document}

% ==================================================================== MASTHEAD
\thispagestyle{empty}
\begin{tikzpicture}[remember picture, overlay]
  \fill[Deep] (current page.north west) rectangle
      ($(current page.north east)+(0,-67mm)$);
  \fill[Teal] ($(current page.north west)+(0,-67mm)$) rectangle
      ($(current page.north east)+(0,-68.4mm)$);
  \foreach \x/\y/\r in {22/12/0.9, 41/26/1.5, 63/9/0.7, 88/31/1.1, 112/17/1.9,
                        138/29/0.8, 159/13/1.3, 178/54/1.0, 196/21/1.6,
                        30/44/0.6, 74/57/1.2, 125/48/0.9, 168/6/0.8, 191/59/0.7,
                        52/18/1.0, 100/60/0.6, 148/20/1.1, 66/49/1.4, 205/45/0.8}
    \fill[Sky, opacity=0.30] ($(current page.north west)+(\x mm,-\y mm)$) circle (\r mm);
\end{tikzpicture}

\vspace*{4mm}
{\sffamily\fontsize{7.6}{9}\selectfont\color{Sky}%
 WEEKLY ROLLUP\;\textbullet\;PhD RESEARCH WATCH\par}
\vspace{2.2mm}
{\sffamily\bfseries\fontsize{27}{29}\selectfont\color{white}%
 Particulate Matter\par}
{\sffamily\fontsize{27}{29}\selectfont\color{Sky}%
 Week in Review\par}
\vspace{5.0mm}
{\sffamily\fontsize{8.4}{10}\selectfont\color{white}%
 Week ending \DIGESTDATE\;\;\textcolor{Teal}{\textbar}\;\;\PERIODSTART{} to \PERIODEND\;\;%
 \textcolor{Teal}{\textbar}\;\;146 records pooled from 7 daily issues\par}
\vspace{18mm}

% -------------------------------------------------------------- metric strip
\noindent
\begin{minipage}[t]{0.190\linewidth}\metric{146}{RECORDS POOLED\\(145 UNIQUE DOIs)}{Deep}\end{minipage}\hfill
\begin{minipage}[t]{0.190\linewidth}\metric{54\%}{ON THE MEASUREMENT\\SIDE (79 RECORDS)}{Teal}\end{minipage}\hfill
\begin{minipage}[t]{0.190\linewidth}\metric{10}{SUBTOPICS\\ACTIVE}{Sky}\end{minipage}\hfill
\begin{minipage}[t]{0.190\linewidth}\metric{30}{POOLED EFFECT\\ESTIMATES}{Amber}\end{minipage}\hfill
\begin{minipage}[t]{0.190\linewidth}\metric{21}{TRIAL REGISTRATIONS\\SURFACED}{Clay}\end{minipage}

\vspace{6mm}

% -------------------------------------------------------------- week in one box
\section*{\sffamily\bfseries\color{Deep}The week in one paragraph}
\vspace{-1.5mm}{\color{Teal}\rule{\linewidth}{1.1pt}}\vspace{2.5mm}

\begin{keybox}
{\sffamily\bfseries\small\color{Deep}This was the week the low-cost-sensor calibration
literature stopped asking whether a correction fits and started asking where it stops
working --- and the same week the epidemiology independently concluded that
source-\ and co-exposure-resolved metrics beat total mass.}\\[1.4mm]
\footnotesize
Twenty-four of the 146 pooled records are low-cost-sensor or calibration papers, and the
strong ones share a single framing. Ginsburg et al.\ tested the US EPA PurpleAir
correction \emph{worldwide} and found it requires regional recalibration; Zhivkov et al.\
quantified calibration-transfer error as a function of \emph{distance} from the reference
node; Chung et al.\ characterised PurpleAir deployment \emph{longevity} across the US
2016--2025; Gäbel \& Hertig examined \emph{recalibration} practice against sensor
ageing. None of these is a new correction equation. All four are about the boundary of an
existing one. On the health side the convergence is independent: Cantuaria et al.\
(\textit{EHP}, \textbf{159{,}769} pooled Danish participants) found that decomposing
PM$_{2.5}$ by source inverts which fraction is harmful relative to UFP and elemental
carbon (total PM$_{2.5}$ \textbf{HR 1.11, 1.05--1.17} per IQR for incident COPD);
Hur et al.\ (\textbf{262{,}288} Kyushu deaths) found the SPM--respiratory-mortality
association separates on pollen load (\textbf{1.021, 1.005--1.037} on high-pollen days
against \textbf{1.003, 0.987--1.018} on low); Sassano et al.\ found household biomass
fuel carried \textbf{HR 1.20 (1.04--1.37)} for kidney-disease mortality with the estimate
splitting on chimney presence. \textbf{For this to be a real convergence rather than a
reading imposed on it, the two literatures would have to be citing each other. They are
not.} Not one of the 146 records validated an instrument against a health endpoint. The
convergence is on the \emph{demand} --- resolved exposure --- and the gap between the two
sides is the finding, not the agreement.
\end{keybox}

\vspace{3mm}

% -------------------------------------------------------------- trend & inflection
\section*{\sffamily\bfseries\color{Deep}Subtopic trend and inflection}
\vspace{-1.5mm}{\color{Teal}\rule{\linewidth}{1.1pt}}\vspace{2.5mm}

\noindent\includegraphics[width=\linewidth]{\FIGDIR/w1_subtopic_trend.png}
\figcap{Per-issue subtopic counts across the week from \texttt{state/metrics.csv}.
\textbf{Read the denominators first}: the daily totals were 31, 25, 13, 10, 27, 26 and 14,
a factor of three between the largest and smallest day, so every count here is a share of
a moving base (the 8 August column exists only because \texttt{metrics.csv} was rebuilt
before this rollup was compiled; a rollup run before its own last daily silently drops
that day). The one genuine inflection is 2 August, when the sensing cluster reached
12 of 31 records --- and roughly half of that spike is \emph{this pipeline}, not the
literature: the Crossref-by-ISSN leg was widened that day and immediately surfaced ten
calibration papers that the PubMed legs structurally cannot see. The 4--5 August trough
(13 and 10 records) is a mid-week Elsevier deposit lull, and the 8 August floor of 14 is
a Saturday. \textbf{None of the three is a signal about the field.} The only movement
that survives normalisation is the steady 32--54\,\% measurement-side share, which held
every single day of the week.}

\footnotesize
\textbf{Emerging this week.} \emph{Calibration transferability as an object of study in
its own right} --- absent from the preceding week, and this week carried by four
independent groups (Ginsburg, Zhivkov, Chung, Gäbel \& Hertig) plus Barkjohn et al.'s
Air Sensor Network Analysis Tool, an R-Shiny QA/QC application that is infrastructure
rather than a result. \emph{Network design as an optimisation problem} also appears for
the first time: Jafari et al.\ jointly optimise bus-mounted mobile sensing against fixed
reference placement, and Bhardwaj \& Subramanian pose identifiability-aware source
apportionment for sparse city-scale networks. Both treat siting as a decision variable
rather than a constraint, which is new in this corpus.

\vspace{1.2mm}
\textbf{Fading.} Neuro and mental-health work collapsed from 22 records in the preceding
week to \textbf{2}, and reproductive/developmental from 24 to 4. \textbf{This is almost
certainly journal-issue timing, not decline.} The preceding week's count was inflated by a
164-record seed backfill on 26 July that included several review clusters; the honest
comparison is 146 records over 7 issues this week against 131 over 6 issues in the
preceding week excluding that backfill --- 20.9 against 21.8 per issue, i.e.\ flat.
Do not read the subtopic deltas against the raw 295.

\vspace{1.2mm}
\textbf{Flat but load-bearing.} Exposure assessment and modelling (32 records) and
mechanistic toxicology (11) held their share through both the peaks and the troughs.
Chemical-transport modelling, satellite retrieval and source apportionment are the
constant substrate of this corpus; they generate few headlines and most of the methods
everything else depends on.

%% LAYOUT RULE (2026-08-02): architecture donut in the whitespace here, not on its own page.
\vspace{3mm}
\noindent\includegraphics[width=\linewidth]{\FIGDIR/f2a_architecture.png}
\figcap{Pooled study architecture over 146 records. Modelling/inventory and measurement
campaigns dominate, which is what an instrumentation-weighted corpus looks like.
\textbf{Reviews and syntheses are a publication-cycle artefact, not new evidence} --- count
them separately before reading the corpus as a body of primary work. \textbf{The
\textit{Trial / intervention} slice is 1 record in 146} --- a fellowship curriculum
evaluation with a pre/post knowledge scale, i.e.\ not a PM intervention at all. Three
further records have intervention-adjacent designs that map elsewhere in this grouping
(a chamber factorial, a secondary analysis of trial cohort data, a filtration-materials
bench study). \textbf{Not one randomised intervention on a PM exposure entered the
corpus in seven days} --- and that is the finding, because 21 such trials are sitting in
the registry watch below.}

\clearpage

% -------------------------------------------------------------- pooled analytics
\section*{\sffamily\bfseries\color{Deep}Pooled analytics}
\vspace{-1.5mm}{\color{Teal}\rule{\linewidth}{1.1pt}}\vspace{3mm}

\noindent\includegraphics[width=\linewidth]{\FIGDIR/f1_subtopics.png}
\figcap{Pooled subtopic distribution, \PERIODSTART{} to \PERIODEND. \textbf{Sensing,
forecasting \& instrumentation (47) and Exposure assessment \& modelling (32) together
account for 79 of 146 records, 54\,\% of the week.} All ten canonical clusters were
active, which is a property of a seven-day window rather than of any single day --- three
clusters were empty on 8 August alone. The clinical tail (Other clinical endpoints 6,
Reproductive \& developmental 4, Neuro 2) is one to two studies per cluster and carries no
weekly trend information.}

\vspace{2mm}
\noindent\includegraphics[width=\linewidth]{\FIGDIR/f2b_endpoint.png}
\figcap{Pooled health endpoint by organ system. The large no-health-endpoint share is the
direct consequence of the 54\,\% measurement-side weighting above: an instrumentation
corpus is mostly records with no human outcome, and that is expected rather than a defect.
\textbf{Eight records this week were carried on metadata alone} --- surfaced with no
abstract in Crossref, Europe PMC or Semantic Scholar --- and they have no endpoint by
construction, not by finding. \textbf{Every bar below 3 records is one or two studies},
not a trend; do not read the tail.}

\clearpage

\noindent\includegraphics[width=\linewidth]{\FIGDIR/f5_forest.png}
\figcap{\textbf{Pooled effect estimates for the week --- read as a caterpillar plot, not a
meta-analysis.} No summary diamond is drawn and none should be. \textbf{Pooling is not
defensible here and no $I^2$ is reported}, for four separate reasons, each sufficient on
its own. (i) \emph{Exposure contrasts differ}: per 10\,$\mu$g/m$^3$ (Orr), per IQR
(Cantuaria), per 10 years of fuel use (Sassano), per binary detection (Tillery), per
hourly mean (Shukla). (ii) \emph{Estimands differ}: HR, RR, OR, IRR and one
percent-averted figure share an axis but not an interpretation. (iii) \emph{Records are
not independent} --- Orr contributes six state-stratified estimates, Cantuaria four
source-resolved ones, Hur three pollen strata and Sassano four ventilation strata; four
papers supply 17 of the 30 points. (iv) \emph{Not all exposures are PM}: Li et al.\ and
Liu et al.\ are explicitly flagged co-pollutant in the corpus, and Shukla and Tillery are
NO$_2$ and nicotine. \textbf{The one subset that is genuinely comparable} is Orr's six
wildfire-PM$_{2.5}$--influenza estimates, all per 10\,$\mu$g/m$^3$ on the same outcome
definition --- and they range from \textbf{0.884 (0.842--0.919)} to \textbf{1.067
(1.056--1.078)} with non-overlapping intervals, which is the week's most important
heterogeneity precisely because it is \emph{within} a single paper and tracks the
outcome-ascertainment boundary, not geography.}

\vspace{2mm}
\noindent\includegraphics[width=\linewidth]{\FIGDIR/f4_heatmap.png}
\figcap{Pooled subtopic $\times$ study-architecture cross-tabulation over the full week.
\textbf{The cells that stay empty are the point.} At $n=146$ --- seven times a daily
corpus --- the \textit{Sensing} and \textit{Exposure assessment} rows still contain no
\textit{Trial / intervention} cell and no cell that couples an instrument to a health
outcome. \textit{Mechanistic toxicology} remains entirely experimental with no bridge to
human biomarker work, the same structural gap this watch recorded on 27 July. A week of
data does not close a cell that a day of data left empty; it establishes that the cell is
empty by construction of the field, not by sampling.}

\clearpage

% -------------------------------------------------------------- groups & citations
\section*{\sffamily\bfseries\color{Deep}Groups, venues and most-cited work}
\vspace{-1.5mm}{\color{Teal}\rule{\linewidth}{1.1pt}}\vspace{2.5mm}

\footnotesize
\textbf{Recurring author and group clusters.} The naive first-author count is
\textbf{misleading and must not be used}: the top entries are Li (5), Wang (4), Zhang (3),
Liu, Zhao, Guo, Yang and Kim (2 each), which are among the most common surnames in the
world and almost certainly denote distinct groups. Two repeats survive inspection.
\textbf{Orr} appears twice in genuinely different roles --- the \textit{EHP} wildfire-smoke
and influenza analysis across six US states, and a \textit{JMIR} feasibility study of an
in-home air-quality monitor in families of children with asthma --- and those two records
are the closest thing this week has to the instrument--endpoint bridge that the heatmap
says is missing. \textbf{Kumpika} appears twice because \textbf{the same paper shipped in
two issues}: DOI \texttt{10.1080/26395940.2026.2616098} was carried on 2 August as
``PMS7003 light-scattering sensors'' and again on 3 August under a broader title.
\textbf{This is a deduplication failure, one record in 146, and it is stated here rather
than quietly corrected} --- the pooled count is 146 records over 145 unique DOIs.
Beyond that, repeat appearances in this corpus are the exception: 146 records, no
institution with more than two, and no evidence of a single programmatic output stream
inflating apparent consistency.

\vspace{2mm}
\textbf{Venue concentration.} \textit{Atmospheric Chemistry and Physics} is the largest
single contributor at \textbf{11 records}, and that is a themed concentration rather than
a research surge --- ACP publishes continuously and the deposits cluster by editorial
batch. After normalising for a label defect described below: \textit{Environmental
Pollution} 8, \textit{Atmospheric Environment} 7, \textit{Research Square} (preprint) 6,
\textit{Atmospheric Measurement Techniques} 5, \textit{Journal of Hazardous Materials} 5,
\textit{Environmental Monitoring and Assessment} 5, \textit{Atmospheric Pollution
Research} 4, \textit{J Air \& Waste Management Association} 4, \textit{Environmental
Health Perspectives} 3. \textbf{Defect worth recording: journal names are stored as free
text and the same venue appears under two spellings} --- \texttt{Atmos Environ}\,(4) and
\texttt{Atmospheric Environment}\,(3), \texttt{Environ Pollut}\,(5) and
\texttt{Environmental Pollution}\,(3), \texttt{Res Sq}\,(2) and \texttt{Research
Square}\,(4). Uncorrected, the register understates venue concentration by roughly a
third at the top of the list.

\vspace{2mm}
\textbf{Most-cited work of the period.} \textbf{Not reportable, and ranking it would be
noise.} Every record in this corpus entered a PubMed or Crossref index within the last
seven days; citation counts are zero or near-zero for all 146 and carry no signal about
importance. Consensus returns citation counts, but for the records it surfaced this week
those counts belong to 2019--2025 papers that were \emph{not} in the entry window and were
rejected for exactly that reason. A defensible most-cited list requires a lag of at least
18 months and will be reported in the yearly issue, not weekly.

\vspace{4mm}

% -------------------------------------------------------------- trial watch
\section*{\sffamily\bfseries\color{Deep}Trial registry watch}
\vspace{-1.5mm}{\color{Teal}\rule{\linewidth}{1.1pt}}\vspace{2.5mm}

{\fontsize{7.4}{9}\selectfont\color{Slate}From \texttt{state/trials.json}: registrations
surfaced inside the window. Daily issues accumulate; the rollup is where they surface.
\textbf{21 registrations} entered state during 2--8 August --- 16 from the 2 August sweep,
1 on 4 August, 3 on 7 August and 1 on 8 August. Note a hard limitation of the connector:
it exposes \emph{no} last-update-date filter, so these are relevance sweeps, not true
window queries, and \texttt{first\_seen} means ``first seen by this pipeline'', not
``registered or updated on that date''.\par}

\vspace{1.5mm}
\footnotesize
\textbf{The structural finding first.} \textbf{Of the 21, at least 13 are sham- or
placebo-controlled portable-air-cleaner or filtration trials.} That is the intervention
literature in this field almost in its entirety, and it has a common design problem:
a HEPA unit in a home typically buys a 30--60\,\% reduction in indoor PM$_{2.5}$ over a
few months, which is a small absolute exposure contrast, and the trials are then powered
on \textbf{surrogate endpoints} --- blood pressure, spirometry, biomarkers, cognitive
scores --- rather than events. Where a hard outcome is attempted, enrolment is far too
small for it.

\vspace{1.5mm}
\textbf{Adequately powered for their stated endpoint (surrogate).}
\textbf{NCT06376994} (JHSPH Center for Clinical Trials, air cleaner vs sham in COPD,
\textbf{n=770}, primary completion 2029-08-31) and \textbf{NCT05874479} (NYU Langone,
active vs sham PAC in NYC public housing, blood pressure, \textbf{n=440}, 2027-09) are the
two that can plausibly detect a several-mmHg or modest spirometric difference.
\textbf{NCT06070428} (Henry Ford, PATHFINDER, PAC at discharge in heart failure,
\textbf{n=400}, 2027-06) is the only one of the three with an event-adjacent population,
and even there the endpoint is not a hard event.

\vspace{1.5mm}
\textbf{Underpowered or endpoint-mismatched --- flagged.} \textbf{NCT06534385}
(Alaska Native Tribal Health Consortium, HEPA purifiers in rural Alaska homes,
\textbf{n=15}) and \textbf{NCT07118189} (Kaiser Permanente, N95 + indoor air cleaner for
wildfire smoke, \textbf{n=20}) cannot support any inferential health claim and should be
read as feasibility work. \textbf{NCT07536178} (UBC, controlled lodgepole-pine woodsmoke
vs filtered air, n=24) is a mechanistic exposure study, correctly small, and is the
cleanest causal design in the list. \textbf{NCT06197477} (Hasselt, school purifiers,
\textbf{n=900}) is the largest trial in state and carries \textbf{four cognitive
co-primary endpoints with no secondaries and no stated multiplicity control} --- with
status \texttt{UNKNOWN} past a December 2025 primary completion, it is now also a
publication-risk item.

\vspace{1.5mm}
\textbf{New on 8 August.} \textbf{NCT05368493} (New York State Psychiatric Institute,
prenatal air pollution and neurodevelopment in Puerto Rican children,
\textbf{n=182}, \texttt{ENROLLING\_BY\_INVITATION}, primary completion 2028-06-01).
\texttt{analyze\_endpoints} returns a \emph{single} primary --- bilateral middle frontal
cortex and corpus callosum volumes, scanned at either 6--36 months or 6--11 years --- and
\textbf{zero secondary endpoints}. Two problems worth recording now rather than at
readout: the registry classifies it \texttt{INTERVENTIONAL} when the only listed
intervention is an MRI task, so it is an observational neuroimaging cohort mislabelled;
and the primary is measured at two non-overlapping developmental windows depending on
enrolment timing, which makes the effective sample for either window well under 182.
Exposure is modelled ambient PM, not measured.

\vspace{2mm}
\textbf{Standing watch list --- results due or overdue.} \textbf{NCT05994937}
(NYU Langone, PAC glycaemic/cardiometabolic biomarkers, n=150) reached primary completion
\textbf{2026-03-16} with \textbf{no results posted} --- five months overdue.
\textbf{NCT06534385} (2026-04-01) and \textbf{NCT07047430} (U. Washington, airport-adjacent
asthma, 2026-06-01) are also past completion. \textbf{NCT07111208} (MARKOPOLO, noise and
UFP, n=180) had primary completion \textbf{2026-07} and is the only UFP-exposure
interventional registration in state. \textbf{NCT07344987} (Peking Union, ILD crossover,
2026-12) and \textbf{NCT07118189} (2026-12-31) close this year.

\clearpage

% -------------------------------------------------------------- gaps & direction
\section*{\sffamily\bfseries\color{Deep}Open gaps}
\vspace{-1.5mm}{\color{Teal}\rule{\linewidth}{1.1pt}}\vspace{2.5mm}

\footnotesize
\begin{itemize}
\item \textbf{Measurement--epidemiology coupling: still zero.} Across 146 records and 24
low-cost-sensor papers, \textbf{not one validated an instrument against a health
endpoint}. The nearest approach is Orr et al.\ (\textit{JMIR}), which measured
feasibility and acceptability of an in-home monitor in families of children with asthma
but reported no exposure--response. The two literatures were adjacent all week and never
touched.

\item \textbf{Ultrafine and speciated metrics.} Five records mention UFP; \textbf{exactly
one measures it as an exposure} (Cantuaria et al., where the local-traffic fraction
dominates for UFP and elemental carbon while the \emph{non}-traffic fraction dominates for
PM$_{2.5}$ --- an inversion with direct siting implications). The rest are a narrative
review, a dosimetry update, a filtration-CFD paper and a mechanistic pathway study.
Against 55 PM$_{2.5}$-only and 18 joint PM$_{2.5}$/PM$_{10}$ records, the deposition
physics remains badly under-served.

\item \textbf{Geographic coverage vs.\ attributable burden.} China 21, North America
$\approx$24, Europe $\approx$10 against \textbf{South Asia $\approx$9 and sub-Saharan
Africa 5}. India and Bangladesh together carry a larger share of global PM-attributable
DALYs than China and North America combined and contributed fewer than ten records in
seven days. The one sub-Saharan primary instrumentation record --- Rehbein et al.\ on
source triangulation from a low-cost network in Ile-Ife --- is exactly the kind of work the
gap needs, and it is one paper.

\item \textbf{Calibration transferability: improving, and now the week's strongest
theme.} Unlike previous weeks, this corpus \emph{does} contain held-out and out-of-region
evaluation: Ginsburg et al.\ (US correction applied worldwide), Zhivkov et al.\
(error as a function of distance from reference), Gomez et al.\ (ML/DL against a T640X
reference), Chung et al.\ (deployment longevity, US 2016--2025). \textbf{What is still
missing is a common benchmark} --- every one of these evaluates on its own network, so the
numbers are not comparable across papers.

\item \textbf{Negative and null results.} Present, which is a good sign. Russo et al.\
report \textbf{protective} point estimates in a spermatogenesis-window analysis
(\textbf{RR 0.71, 0.55--0.93} and 0.73, 0.50--1.05) that are more likely selection or
window misspecification than a real benefit; Orr et al.\ report a protective estimate in
Washington state (\textbf{0.884, 0.842--0.919}) that the paper itself cannot explain
except by outcome misclassification; Cantuaria reports two null source-resolved HRs
(1.02, 1.00--1.05 and 1.05, 0.99--1.11). \textbf{A corpus with no nulls is a
reporting-bias signal; this one has them, and they are informative rather than
decorative.}
\end{itemize}

\vspace{3mm}
\begin{keybox}
{\sffamily\bfseries\footnotesize\color{Deep}Where this is heading --- judgement, clearly
marked as such}\\[1.2mm]
\footnotesize
\textbf{The methodological demand the field is converging on is a resolved exposure
record, and the convergence is arriving from two directions that are not talking to each
other.} From epidemiology: source apportionment (Cantuaria), co-pollutant interaction
(Hur), microenvironment (Sassano, Tillery) and sub-daily lag structure (Shukla) each
independently showed this week that a scalar mass concentration is not a sufficient
statistic for the health-relevant exposure. From instrumentation: four groups
simultaneously stopped publishing new correction equations and started publishing the
\emph{failure boundaries} of existing ones. Both are asking for the same object. Neither
is building it. Confidence: moderate-to-high on the description, low on any claim that
this week caused it --- seven days is not enough to establish a turn, and 2 August's
calibration cluster was partly an artefact of widening this pipeline's own Crossref leg.

\vspace{1.4mm}
\textbf{What would change the picture.} A single paper that deploys a calibrated low-cost
network and reports an exposure--response against a clinical endpoint would collapse the
empty heatmap cell that has now survived two weeks of accumulation. Nothing in the 21
registered trials will supply it: they are almost all filtration interventions with
surrogate endpoints, and the two that could (NCT06376994, n=770; NCT05874479, n=440) read
out in 2027 and 2029. If the coupling comes, it will come from a cohort retrofitting a
dense sensor network onto existing outcome ascertainment, not from a trial.

\vspace{1.4mm}
\textbf{What a distributed low-cost PM sensing programme should do differently.} Three
concrete things this week's corpus argues for. \textbf{(1)} Report calibration
performance as a function of distance and time from the co-location reference, not as a
single held-out R$^2$ --- Zhivkov et al.\ give the framing and Chung et al.\ give the
longevity data to do it. \textbf{(2)} Instrument the co-exposure, not just the mass:
Hur's pollen interaction and Cantuaria's source split both say the node needs a second
channel (fluorescence, coarse-mode, or a traffic-proxy gas) before its mass output is
epidemiologically interpretable. \textbf{(3)} Resolve below 24 hours and keep it ---
Shukla localises the asthma signal at a 3--6 hour lag, which no filter-based reference
method can see and which is the one regime where a low-cost network is not merely cheaper
but strictly better. \textbf{The week was informative, and unusually so on the
instrumentation axis specifically.}
\end{keybox}

\clearpage

% ==================================================================== DIGEST
\section*{\sffamily\bfseries\color{Deep}Paper-level digest --- the week by subtopic}
\vspace{-1.5mm}{\color{Teal}\rule{\linewidth}{1.1pt}}\vspace{2mm}

\footnotesize
Every record of the period, grouped by subtopic. \textbf{Two tiers, and the difference
matters.} Entries carrying a full critical summary are lifted \emph{verbatim} from the
archived daily issue that first published them --- they were read, and the stated
weaknesses are assessed weaknesses. The remainder appear as precis or metadata-only
entries in the lighter style: those records were surfaced and classified but
\textbf{not read}, either because no abstract was available in Crossref, Europe PMC or
Semantic Scholar (8 records this week) or because they entered on metadata alone.
\textbf{Absence of criticism in the lighter tier means the record is unassessed, not
that it is sound.} Do not cite a precis entry as though it had been evaluated.

\vspace{3mm}

\input{digest_body}

\clearpage

% ==================================================================== TABLE
\section*{\sffamily\bfseries\color{Deep}Pooled corpus register}
\vspace{-1.5mm}{\color{Teal}\rule{\linewidth}{1.1pt}}\vspace{2.5mm}

{\fontsize{6.7}{8.4}\selectfont
\setlength{\tabcolsep}{2.4pt}
\begin{longtable}{@{}>{\RaggedRight}p{25mm}>{\RaggedRight}p{41mm}>{\RaggedRight}p{27mm}%
>{\RaggedRight}p{23mm}>{\RaggedRight}p{18mm}>{\RaggedRight}p{15mm}@{}}
\toprule
\textsf{\textbf{First author}} & \textsf{\textbf{Focus}} & \textsf{\textbf{Journal}} &
\textsf{\textbf{Design}} & \textsf{\textbf{Metric}} & \textsf{\textbf{Region}}\\
\midrule
\endfirsthead
\toprule
\textsf{\textbf{First author}} & \textsf{\textbf{Focus}} & \textsf{\textbf{Journal}} &
\textsf{\textbf{Design}} & \textsf{\textbf{Metric}} & \textsf{\textbf{Region}}\\
\midrule
\endhead
\multicolumn{6}{@{}l@{}}{\cellcolor{Mist}\textsf{\textbf{\textcolor{Deep}{Neuro / mental health}}}}\\[0.6mm]
\href{https://doi.org/10.1001/jamanetworkopen.2026.25015}{Dales} & Wildfire-sourced PM$_{2.5}$ and ED visits for migraine and other primary headache & \textit{JAMA Netw Open} & Case-crossover & PM$_{2.5}$ (wildfire) & Canada\\
\href{https://doi.org/10.1093/ageing/afag215}{Glover} & Umbrella meta-review of urban environmental exposures and cognitive health & \textit{Age Ageing} & Umbrella review / meta-analysis & PM$_{2.5}$ & Global\\
\href{https://doi.org/10.1016/j.nsa.2026.107019}{Goossens} & In utero air pollution exposure and pre-/postnatal brain development & \textit{Neurosci Appl} & Narrative review & PM$_{2.5}$ / PM$_{10}$ / PM$_{\mathrm{coarse}}$ & Global\\
\href{https://doi.org/10.1016/j.isci.2026.116617}{Ma} & PM$_{2.5}$-induced neurotoxicity and depression: exposure risks and mechanisms & \textit{iScience} & Narrative review & PM$_{2.5}$ + components & Global\\
\href{https://doi.org/10.3760/cma.j.cn112150-20250804-00752}{Pan} & Maternal PM exposure and offspring neurodevelopment: effects and mechanisms & \textit{Zhonghua Yu Fang Yi Xue Za Zhi} & Narrative review & PM (bulk) & China\\
\addlinespace[0.9mm]
\multicolumn{6}{@{}l@{}}{\cellcolor{Mist}\textsf{\textbf{\textcolor{Teal}{Cardiovascular \& metabolic}}}}\\[0.6mm]
\href{https://doi.org/10.3760/cma.j.cn112338-20251120-00843}{Huang} & PM$_{2.5}$ and dyslipidaemia in rural Chinese children; dietary effect modification & \textit{Zhonghua Liu Xing Bing Xue Za Zhi} & Cross-sectional (multistage) & PM$_{2.5}$ & China\\
\href{https://doi.org/10.32394/pe/218966}{Kazmierska} & Air pollution and cardiovascular disease: review of current data & \textit{Przegl Epidemiol} & Narrative review & PM$_{2.5}$ / PM$_{10}$ & Global\\
\href{https://doi.org/10.1007/s00125-026-06814-2}{Xi} & Air pollution, built environment and temperature vs. prediabetes phenotypes & \textit{Diabetologia} & Prospective cohort & PM$_{2.5}$ & Germany\\
\href{https://doi.org/10.1161/JAHA.125.046917}{Yang} & Short-term PM$_{2.5}$ and hospitalisation/ED use in atrial fibrillation & \textit{J Am Heart Assoc} & Case-crossover & PM$_{2.5}$ (wildfire-influenced) & USA\\
\addlinespace[0.9mm]
\multicolumn{6}{@{}l@{}}{\cellcolor{Mist}\textsf{\textbf{\textcolor{Sky}{Reproductive \& developmental}}}}\\[0.6mm]
\href{https://doi.org/10.1016/j.jes.2025.10.043}{Chen} & Prenatal air pollution, fetal growth trajectory phenotypes and child neurodevelopment & \textit{J Environ Sci (China)} & Prospective cohort & PM$_{2.5}$ / PM$_{10}$ & China\\
\href{https://doi.org/10.1016/j.ecoenv.2026.120566}{Huang} & Joint prenatal bisphenol + air-pollutant mixture exposure and preschool IQ & \textit{Ecotoxicol Environ Saf} & Prospective cohort (mixtures) & PM$_{2.5}$ / PM$_{10}$ & China\\
\href{https://doi.org/10.1007/s40572-026-00552-8}{Peterson} & Ambient air pollution, PAHs and pubertal development: critical review & \textit{Curr Environ Health Rep} & Critical review & PM$_{2.5}$ components & Global\\
\href{https://doi.org/10.1016/j.envres.2026.125313}{Xie} & Long-term PM$_{2.5}$ and timing of pubertal development in boys and girls & \textit{Environ Res} & Prospective cohort & PM$_{2.5}$ & China\\
\addlinespace[0.9mm]
\multicolumn{6}{@{}l@{}}{\cellcolor{Mist}\textsf{\textbf{\textcolor{Sage}{Respiratory \& allergic}}}}\\[0.6mm]
\href{https://doi.org/10.1007/s00484-026-03289-z}{Bayar Muluk} & Climate change and otolaryngology: airway disease impacts and emission reduction & \textit{Int J Biometeorol} & Narrative review & PM$_{2.5}$ & Global\\
\href{https://doi.org/10.4103/mgr.MEDGASRES-D-25-00164}{Liao} & Short-term air pollution, respiratory symptoms and lung function in asthma & \textit{Med Gas Res} & Panel study & PM$_{2.5}$ / PM$_{10}$ & China\\
\href{https://doi.org/10.3390/epidemiologia7040097}{Sanchez de Prada} & Ambient PM and asthma-coded admissions in children 0-3 y, Spain & \textit{Epidemiologia} & Ecological panel & PM$_{2.5}$ / PM$_{10}$ & Spain\\
\href{https://doi.org/10.1016/j.jhazmat.2026.143076}{Yang} & Allergenic vs. toxic risk of airborne seed hairs compared with ambient PM$_{2.5}$ & \textit{J Hazard Mater} & Field measurement & PM$_{2.5}$ + bioaerosol & China\\
\addlinespace[0.9mm]
\multicolumn{6}{@{}l@{}}{\cellcolor{Mist}\textsf{\textbf{\textcolor{Violet}{Mechanistic toxicology}}}}\\[0.6mm]
\href{https://doi.org/10.1016/j.jes.2025.12.074}{Ding} & PM$_{2.5}$ promotes vascular calcification via miR-27b-5p/HES1 signalling & \textit{J Environ Sci (China)} & In vitro + organoid + murine & PM$_{2.5}$ & China\\
\href{https://doi.org/10.1007/s00360-026-01698-5}{Maraschi} & Osmoregulatory disruption and metal bioaccumulation from settleable atmospheric PM & \textit{J Comp Physiol B} & Controlled exposure (ecotox) & Settleable APM & Brazil\\
\href{https://doi.org/10.1016/j.jes.2025.11.018}{Wu} & Macrophage NLRP3 inflammasome in real-ambient PM-induced hepatic glucose dysregulation & \textit{J Environ Sci (China)} & Real-ambient murine + co-culture & Real-ambient PM & China\\
\addlinespace[0.9mm]
\multicolumn{6}{@{}l@{}}{\cellcolor{Mist}\textsf{\textbf{\textcolor{Clay}{Occupational \& indoor}}}}\\[0.6mm]
\href{https://doi.org/10.3389/fmedt.2026.1796828}{Chen} & Nested suction instrument to cut intraoperative surgical-smoke PM$_{2.5}$ & \textit{Front Med Technol} & Bench / device experiment & PM$_{2.5}$ (surgical smoke) & China\\
\href{https://doi.org/10.1093/annweh/wxag055}{Cole} & Determinants of respirable crystalline silica exposure in Australian tunnelling & \textit{Ann Work Expo Health} & Observational exposure study & Respirable dust / RCS & Australia\\
\href{https://doi.org/10.1016/j.jhazmat.2026.143081}{Ding} & Stage-resolved hair Zn-isotope dynamics as a firefighting combustion-exposure tracer & \textit{J Hazard Mater} & Repeated-measures biomarker & Combustion-derived PM & China\\
\href{https://doi.org/10.3390/jox16040119}{Ramsden} & Xenobiotic hazards in aircraft cabin air & \textit{J Xenobiot} & Review + risk economics & Ultrafine particles & Global\\
\addlinespace[0.9mm]
\multicolumn{6}{@{}l@{}}{\cellcolor{Mist}\textsf{\textbf{\textcolor{Amber}{Sensing, forecasting \& instrumentation}}}}\\[0.6mm]
\href{https://doi.org/10.1016/j.jenvman.2026.130515}{Shi} & Adaptive spatiotemporal graph transformer for multi-site multi-step PM$_{2.5}$ forecasting & \textit{J Environ Manage} & Deep-learning model & PM$_{2.5}$ & China / India\\
\href{https://doi.org/10.3390/bios16070380}{Tsuchiya} & Seasonal and indoor-outdoor characteristics of fluorescent aerosol particles in offices & \textit{Biosensors} & Real-time sensor campaign & Bioaerosol / coarse & Japan\\
\addlinespace[0.9mm]
\multicolumn{6}{@{}l@{}}{\cellcolor{Mist}\textsf{\textbf{\textcolor{Amber}{Exposure assessment \& modelling}}}}\\[0.6mm]
\href{https://doi.org/10.1029/2025GH001621}{Solorzano-Araque} & Multiscale exposure modelling: LOTOS-EUROS + mobility surveys + particle toxicity & \textit{GeoHealth} & CTM + mobility integration & PM$_{2.5}$ / PM$_{10}$ & Colombia\\
\addlinespace[0.9mm]
\multicolumn{6}{@{}l@{}}{\cellcolor{Mist}\textsf{\textbf{\textcolor{Coral}{Burden, policy \& mitigation}}}}\\[0.6mm]
\href{https://doi.org/10.1002/hsr2.72886}{Karamzad} & Ischaemic heart disease burden in MENA 1990-2021; PM attributable fraction 28.2\% & \textit{Health Sci Rep} & GBD secondary analysis & Ambient PM & MENA\\
\href{https://doi.org/10.1371/journal.pone.0347989}{Rees} & HEPA filtration to reduce respiratory infection in care homes & \textit{PLoS One} & RCT process evaluation & Filtered PM & UK\\
\href{https://doi.org/10.1016/j.jes.2026.01.009}{Shan} & Policy-driven emission reductions and post-pandemic rebound, high-resolution inventories & \textit{J Environ Sci (China)} & Emission inventory & PM emissions & China\\
\addlinespace[0.9mm]
\multicolumn{6}{@{}l@{}}{\cellcolor{Mist}\textsf{\textbf{\textcolor{Slate}{Other clinical endpoints}}}}\\[0.6mm]
\href{https://doi.org/10.3390/dj14070440}{Lupita} & PM$_{2.5}$, preventive-care capacity and caries/edentulism burden across the EU & \textit{Dent J} & Ecological panel & PM$_{2.5}$ & EU-27\\
\href{https://doi.org/10.1016/j.jdin.2026.06.008}{Park} & Long-term PM exposure and risk of androgenetic alopecia & \textit{JAAD Int} & Nationwide cohort & PM$_{2.5}$ / PM$_{10}$ & South Korea\\
\href{https://doi.org/10.3171/2026.3.PEDS2615}{Sun} & Environmental exposure and area deprivation vs. paediatric hydrocephalus outcomes & \textit{J Neurosurg Pediatr} & Registry cohort + geospatial & PM$_{2.5}$ (satellite-derived) & USA\\
\addlinespace[0.9mm]
\bottomrule
\end{longtable}}


\vspace{4mm}
\begin{keybox}
{\sffamily\bfseries\footnotesize\color{Deep}Provenance and method}\\[1.2mm]
{\fontsize{7.2}{8.8}\selectfont
This rollup is an \textbf{aggregation only}. It pools \texttt{state/corpus/*.json},
\texttt{state/metrics.csv} and \texttt{state/trials.json} over
\PERIODSTART{}--\PERIODEND{} via \texttt{PMRW\_START}/\texttt{PMRW\_END} and
\textbf{re-queries no API}, so every record was screened and deduplicated in the daily
issue that first carried it. Per-record provenance, source-query windows, exclusions and
known metadata defects are documented in those daily issues and are not restated here.
\textbf{Daily issues covered: 2, 3, 4, 5, 6, 7 and 8 August 2026 --- seven of seven, no
missed runs and no gap in the entry-date coverage.}
\textbf{Known defects carried into this rollup, stated rather than hidden:} (i) one
cross-issue duplicate, DOI \texttt{10.1080/26395940.2026.2616098}, shipped on both 2 and
3 August, so 146 records span 145 unique DOIs; (ii) journal names are free text and three
venues appear under two spellings each, which understates venue concentration in the
register; (iii) 8 records carry no abstract from any checked source and are unassessed;
(iv) several 2--7 August records were Consensus or Crossref backfills with creation dates
outside the window and are excluded from trend claims but included in the pooled counts.
Effect estimates are transcribed from abstracts and use heterogeneous exposure contrasts
--- see the forest-plot caption; \textbf{no meta-analytic summary is computed and none
should be}. Trend and inflection statements are judgements over a seven-day window and are
not statistical tests. Content derived from publisher metadata; abstracts remain the
copyright of their respective publishers.\par}
\end{keybox}

\end{document}
